Describe the equivalence between solvency tests and acceptability of future own
funds.

Let \( A_t \) denote the assets  and \( L_t \) the liabilities of an insurance company such as they're presented on a balance sheet.
Then the own founds/net asset values are \( E_t = A_t - L_t \).
The quantities at time \(t=0\) are known, the quantities at time \(t=1\) are random variables on a given probability space 
\((\Omega, \mathcal{F}, \mathbb{P})\).
The increment of the net asset value is \(\Delta E_1 := E_1 - E_0\).

Solvency test
For a given regulatory monetary risk measure \(\rho\), the company is solvent if
\[\rho(\Delta E_1) \leq E_0 \Leftrightarrow \rho(E_1) \leq 0 \Leftrightarrow E_1 \in A_{\rho}\]

Characterize in more detail the specific solvency tests based on V@R and AV@R
and their relationship to the probability of solvency.

Standard examples are the monetary risk measures Value at Risk (V@R) and Average Value at Risk (AV@R) at some level \(\alpha \in (0,1)\):
\[V@R_{\alpha}(X) := \inf \{ x \in \mathbb{R} : P(X+x<0) \leq \alpha \}\]
\[AV@R_{\alpha}(X) := \frac{1}{\alpha} \int_0^{\alpha} V@R_{\beta}(X) d\beta\]

Regulatory standards
- Solvency II \quad V@R at level \(\alpha = 0.5\%\)
- Swiss solvency test \quad AV@R at level \(\alpha = 1\%\)
- Basel III \quad AV@R with level \(\alpha = 2.5\%\)

Value at Risk and Average Value at Risk
- \(V@R_{\alpha}(\Delta E_1) \leq E_0 \quad \Leftrightarrow \quad P(A_1 \geq L_1) \geq 1-\alpha\)
- \(AV@R_{\alpha}(\Delta E_1) \leq E_0 \quad \Leftrightarrow \quad P(A_1 \geq L_1) \geq 1-\alpha\)


Let \( λ ∈ [0, 1] \) be a recovery fraction, what inclusions hold for the events \( \{A_1 \leq \lambda L_1\}\)
for different \( \lambda \)'s. Formalize the largest lower bound of the probability of these set
under solvency constraints with respect to V@R and AV@R, respectively.

We denote by \(\mathcal{X}\) a set of positive random variables on some nonatomic probability space \((\Omega, \mathcal{F}, P)\). 
We assume that \(\mathcal{X}\) contains all positive discrete random variables and fix \(\alpha \in (0,1)\).
For all \(\lambda \in (0,1)\) we have

\[1 - \alpha = \inf \{ P(A \geq \lambda L) : A, L \in \mathcal{X}, AV@R_{\alpha}(A-L) \leq 0 \}\]
\[= \inf \{ P(A \geq \lambda L) : A, L \in \mathcal{X}, V@R_{\alpha}(A-L) \leq 0 \}.\]


How can this regulatory problem be resolved? Argue in detail.

NOT FOUND

Illustrate the role of guard rails in a simple toy model of the behavior of equity
holders.

Probability space \(\Omega = \{g, b\}\) with \(P(b) = \frac{a}{2}\) with \(a \approx 0\), say \(a = 0.5\%\) or \(a = 1\%\).
Liabilities
\[L_1(\omega) = \begin{cases} 1 \quad \text{if } \omega = g, \\ 100 \quad \text{if } \omega = b. \end{cases} \]
The company can manage its assets by engaging in a stylized financial contract with zero initial cost transferring dollars from the good state to the bad state.
More specifically, we assume that the company can choose one of the following asset profiles at time 1:
\[A_1^k(\omega) = \begin{cases} 101-k \quad \text{if } \omega = g, \\ k \quad \text{if } \omega = b, \end{cases} \quad \text{with } k \in [0, 100]. \]
Hedging its liabilities completely would require the company to choose \(k = 100\).

For any \(k \in [0, 100]\), the company's net asset value is given by
\[E_1^k(\omega) = \begin{cases} 100-k \quad \text{if } \omega = g, \\ k-100 \quad \text{if } \omega = b. \end{cases} \]

Due to limited liability, the corresponding shareholder value is
\[\max\{E_1^k(\omega), 0\} = \begin{cases} 100-k \quad \text{if } \omega = g, \\ 0 \quad \text{if } \omega = b. \end{cases} \]

Hence, the choice \(k=0\) is optimal from the perspective of shareholders — corresponding to no recovery in the bad state.

The company is solvent independent of \(k\):
\[V@R_{\alpha}(E_1^k) = k-100 \leq 0, \quad AV@R_{\alpha}(E_1^k) = \frac{1}{\alpha} \left( \frac{\alpha}{2}(100-k) + \frac{\alpha}{2}(k-100) \right) = 0.\]

What are general recovery risk measures? What is RecV@R/RecAV@R?

The Recovery Value at Risk with increasing level function \(\gamma: [0,1] \rightarrow [0,1)\) is defined by

\[RecV@R_{\gamma}(X, Y) = \sup_{\lambda \in [0,1]} V@R_{\gamma(\lambda)}(X + (1-\lambda)Y).\]

In applications, \(X\) plays the role of the net asset value \(E_1\) or the increment of the net asset value \(\Delta E_1\), and \(Y\) 
refers to liabilities \(L_1\).

For every \(\lambda \in [0,1]\) consider a map \(\rho_{\lambda}: \mathcal{X} \rightarrow \mathbb{R} \cup \{\infty\}\) and assume that 
\(\rho_{\lambda_1} \geq \rho_{\lambda_2}\) whenever \(\lambda_1 \leq \lambda_2\). 
The recovery risk measure \(Rec\rho: \mathcal{X} \times \mathcal{X} \rightarrow \mathbb{R} \cup \{\infty\}\) is defined by 
\(Rec\rho(X, Y) := \sup_{\lambda \in [0,1]} \rho_{\lambda}(X + (1-\lambda)Y)\).

Another example of a general recovery risk measure is recovery average value at risk that dominates \(RecV@R\) and which is 
cash invariant in its first component, monotone, convex, subadditive, positively homogeneous, star shaped in its first component, 
and normalized.

The Recovery Average Value at Risk \(RecAV@R_{\gamma}: L^1 \times L^1 \rightarrow \mathbb{R} \cup \{\infty\}\) with increasing level function \(\gamma: [0,1] \rightarrow [0,1)\) is defined by
\[RecAV@R_{\gamma}(X, Y) := \sup_{\lambda \in [0,1]} AV@R_{\gamma(\lambda)}(X + (1-\lambda)Y).\]


Formulate the Markowitz problem with RecAV@R. How can this be reduced to
an LP which can be solved by standard algorithms like the simplex algorithm

Risk measures are an important instrument to limit downside risk in portfolio optimization problems.
- This idea is related to the classical Markowitz problem (1952) in which standard deviation quantifies the risk. Efficient frontiers characterize 
  the best tradeoffs between return and risk.
- Standard deviation is not a sensible risk measure outside elliptical distributions. In realistic applications, 
  it must be replaced by better downside risk measures.

Recovery risk measures may successfully be applied to portfolio optimization in practice.
For \(RecAV@R\) with suitable recovery functions the characterization of the efficient frontier may, on the basis of a suitable minimax theorem, 
be reduced to the minimization of a linear function on a convex polyhedron.

The Markowitz problem is then:
For a given level function \(\gamma: [0,1] \rightarrow [0,1]\) and for given \(a \in \mathbb{R}\), we consider the following problem
\[\min_{\mathbf{x} \in \Delta^K} RecAV@R_{\gamma}(E_1(\mathbf{x}), L_1)\]
\[\text{s.t. } E(A_1(\mathbf{x})) \geq a.\]

State how the Markowitz problem can be reduced to
an LP which can be solved by standard algorithms like the simplex algorithm



The net asset value of the company equals

\[E_1(\mathbf{x}) := A_1(\mathbf{x}) - L_1 = b \left( 1 + \sum_{k=1}^K X^k R^k - Z \right).\]
where \( x^k > 0\) is the fraction of total budget \( b > 0 \) invested into product \( k \).
We define
\[\Delta^K := \left\{ \mathbf{x} \in \mathbb{R}_+^K; \sum_{k=1}^K x^k = 1 \right\}\] 
\[\Delta^K(\mu) := \left\{ \mathbf{x} \in \Delta^K; \sum_{k=1}^K x^k E(R^k) \geq \mu \right\}.\]

The Markowitz problem then requires us to compute and optimize
\[RecAV@R_{\gamma} \left( b \cdot \left[ 1 + \sum_{k=1}^K x^k R^k - Z \right], bZ \right) = -b + b \cdot RecAV@R_{\gamma} \left( \sum_{k=1}^K x^k R^k - Z, Z \right).\]
We focus on the special case of piecewise-constant recovery function. In this case, \(RecAV@R\) is a maximum of finitely many terms:
\[RecAV@R_{\gamma} \left( \sum_{k=1}^K x^k R^k - Z, Z \right) = \max_{i=1, \dots, n+1} AV@R_{\alpha_i} \left( \sum_{k=1}^K x^k R^k - r_i Z \right).\]

- For convenience, for \(i = 1, \dots, n+1\) we define auxiliary functions for the OCE representation of \(AV@R\):
\[\Psi^i(\mathbf{x}, v) = \frac{1}{\alpha_i} \cdot E \left[ \left( v - \sum_{k=1}^K x^k R^k - r_i Z \right)^+ \right] - v.\]
Minimizing over \(v \in \mathbb{R}\) gives \(AV@R_{\alpha_i} \left( \sum_{k=1}^K x^k R^k - r_i Z \right)\).
- This allows us to write
\[RecAV@R_{\gamma} \left( \sum_{k=1}^K x^k R^k - Z, Z \right) = \max_{i=1, \dots, n+1} \min_{v \in \mathbb{R}} \Psi^i(\mathbf{x}, v).\]
- The original problem can thus be rewritten:
\[\min_{\mathbf{x} \in \Delta^K(\mu)} \max_{i=1, \dots, n+1} \min_{v \in \mathbb{R}} \Psi^i(\mathbf{x}, v).\]

The efficient frontier (balancing downside risks with returns) can then be found via the following theorem:

For every \(\mathbf{x} \in \Delta^K\) the following minimax equality holds:
\[\max_{i=1, \dots, n+1} \min_{v \in \mathbb{R}} \Psi^i(\mathbf{x}, v) = \min_{v \in \mathbb{R}^{n+1}} \max_{i=1, \dots, n+1} \Psi^i(\mathbf{x}, v^i).\]

In particular, problem (3) can be equivalently written as
\[\min_{\mathbf{x} \in \Delta^K(\mu)} \min_{v \in \mathbb{R}^{n+1}} \max_{i=1, \dots, n+1} \Psi^i(\mathbf{x}, v^i).\]

The problem of determining the portfolios with miminal risk for a fixed expected target return \(\mu \in \mathbb{R}\) can be equivalently expressed as
\[\min_{(\mathbf{x}, v, T) \in \Delta^K(\mu) \times \mathbb{R}_+^{n+1} \times \mathbb{R}} \{ T; \Psi^i(\mathbf{x}, v^i) \leq T, i = 1, \dots, n+1 \}.\]
This can be solved via Linear Programming (LP).


Describe types of Knightian uncertainty that could additionally be modeled and
for which solutions can still be obtained.

In the presence of Knightian uncertainty, we assume that the statistical probability measure is not known. Instead we take a worst-case approach with uncertainty set \(\mathcal{M} \subset \mathcal{P}\).

Definition 4.4.1

Let \(\gamma: [0,1] \rightarrow [0,1]\) be an increasing function and \(\mathcal{M} \subset \mathcal{P}\). 
The corresponding Worst-Case Recovery Average Value at Risk is defined by
\[RecAV@R_{\gamma}^{\mathcal{M}}(X, Y) := \sup_{P \in \mathcal{M}} RecAV@R_{\gamma}^P(X, Y).\]

Robust problem

We are interested in the following robust problem:
\[\min_{\mathbf{x} \in \Delta^K} RecAV@R_{\gamma}^{\mathcal{M}}(E_1(\mathbf{x}), L_1)\]
\[\text{s.t. } \inf_{P \in \mathcal{M}} E_P(A_1(\mathbf{x})) \geq a.\]

This in turn can be used to solve for Mixture Uncertainty or Box Uncertainty via a combination of Monte-Carlo and LP methods.
Mixture uncertainty can be succinctly stated as 
\(\mathcal{M}_{\text{mix}} = \{ \sum_{j=1}^m \lambda_j \mathbb{P}_j : \mathbf{\lambda} \in \Delta^m \}\).


